\section{Auditoria das supostas citações literais}
A tabela preserva a identificação e a grafia atribuídas às fontes no material inicial. Não reproduz as frases como citações autênticas. O diagnóstico de não autenticação indica insuficiência de comprovação nesta coleta, não demonstração de inexistência em qualquer edição. A matriz CSV conserva uma âncora curta para buscas futuras.
\small\begin{longtable}{p{0.09\textwidth}p{0.27\textwidth}p{0.19\textwidth}p{0.29\textwidth}}
\caption{Os 22 alvos da série enviada: grafias alegadas, sem autenticação das frases.}\\
\toprule Ano & Fonte alegada & Grafia alegada & Falta documental \\\midrule\endfirsthead
\toprule Ano & Fonte alegada & Grafia alegada & Falta documental \\\midrule\endhead
1902 & Canaã & Ah! Ah! Ah! & Sem edição e página. \\
1904 & Gazeta de Notícias, 25/07 & Ah! Ah! Ah! & Sem página e autoria da crônica. \\
1911 & Triste fim de Policarpo Quaresma & Ah! ah! & Sem edição e página; separar folhetim de 1911 e livro de 1915. \\
1912 & O Malho, 12/10 & Quá! Quá! Quá! & Sem página e número. \\
1928 & Macunaíma & Ah! ah! ah! & Sem página e edição. \\
1925 & Correio da Manhã, 14/02 & Ah! ah! ah! & Sem página e autoria. \\
1934 & São Bernardo & Ah, ah! & Sem página e edição. \\
1935 & Fon-Fon!, 08/06 & Quá! quá! quá! & Sem página. \\
1946 & Sagarana & He! he! he! & Sem conto, página e edição. \\
1944 & Diário de Notícias, 18/03 & (Risos na assistência) & Sem página e identificação do ministro. \\
1953 & Almanaque do Pato Donald & Quá-quá-quá! & Sem número, história e página. \\
1955 & Tribuna da Imprensa & He-he-he & Sem dia, página e título da coluna. \\
1967 & Quarup & Ha, ha, ha & Sem página e edição. \\
1969 & O Pasquim, n. 12 & (hahaha) & Sem página e entrevista. \\
1973 & As meninas & Ha-ha-ha! & Sem página e edição. \\
1975 & Revista do Cebolinha & Hi-hi-hi! & Sem número, história e página. \\
1982 & Feliz ano velho & Ha ha ha & Sem página e edição. \\
1986 & Folha de S.Paulo, Painel, 14/11 & Rs. & Sem página. \\
1993 & Veja, Páginas Amarelas, 07/05 & (Risos) & Entrevistado não identificado. \\
1997 & Folhateen, junho & rsrsrs / hahaha & Sem dia, página e reportagem. \\
2005 & Os segredos do Orkut & kakakaka & Sem autor, editora e edição identificados. \\
2007 & Capricho, setembro & KKKKKKKKKK & Sem número, página e reportagem. \\
\bottomrule\end{longtable}\normalsize

\section{Cobertura anual e distinção entre lacuna e ausência}
Os anos abaixo pertencem ao núcleo 1900--2009. A coluna de atestações conta documentos incluídos no inventário, não todos os arquivos baixados nem o universo publicado. O símbolo -- significa que não há atestação confirmada naquele ano nesta coleta; não representa frequência histórica zero. As extensões de 1872 a 2026 constam no arquivo público.
\small\begin{longtable}{rrp{0.65\textwidth}}
\caption{Cobertura anual do núcleo 1900--2009. Não é série de prevalência.}\\
\toprule Ano & Documentos & Interpretação \\\midrule\endfirsthead
\toprule Ano & Documentos & Interpretação \\\midrule\endhead
1900 & 1 & Atestação confirmada \cite{anual1900}. \\
1901 & 1 & Atestação confirmada \cite{queiros1901}. \\
1902 & 1 & Atestação confirmada \cite{anual1902}. \\
1903 & 1 & Atestação confirmada \cite{anual1903}. \\
1904 & 1 & Atestação confirmada \cite{anual1904}. \\
1905 & 1 & Atestação confirmada \cite{anual1905}. \\
1906 & 1 & Atestação confirmada \cite{anual1906}. \\
1907 & 1 & Atestação confirmada \cite{anual1907}. \\
1908 & 1 & Atestação confirmada \cite{anual1908}. \\
1909 & 1 & Atestação confirmada \cite{anual1909}. \\
1910 & 1 & Atestação confirmada \cite{anual1910}. \\
1911 & 1 & Atestação confirmada \cite{anual1911}. \\
1912 & 1 & Atestação confirmada \cite{anual1912}. \\
1913 & 1 & Atestação confirmada \cite{anual1913}. \\
1914 & 1 & Atestação confirmada \cite{anual1914}. \\
1915 & 1 & Atestação confirmada \cite{anual1915}. \\
1916 & 1 & Atestação confirmada \cite{anual1916}. \\
1917 & 1 & Atestação confirmada \cite{almada1917}. \\
1918 & 1 & Atestação confirmada \cite{anual1918}. \\
1919 & 1 & Atestação confirmada \cite{lobato1919}. \\
1920 & 1 & Atestação confirmada \cite{anual1920}. \\
1921 & 1 & Atestação confirmada \cite{anual1921}. \\
1922 & 1 & Atestação confirmada \cite{anual1922}. \\
1923 & 1 & Atestação confirmada \cite{anual1923}. \\
1924 & 1 & Atestação confirmada \cite{anual1924}. \\
1925 & 1 & Atestação confirmada \cite{portugal1925}. \\
1926 & 1 & Atestação confirmada \cite{anual1926}. \\
1927 & 1 & Atestação confirmada \cite{anual1927}. \\
1928 & 1 & Atestação confirmada \cite{diarionacional1928}. \\
1929 & 1 & Atestação confirmada \cite{anual1929}. \\
1930 & 1 & Atestação confirmada \cite{anual1930}. \\
1931 & 1 & Atestação confirmada \cite{anual1931}. \\
1932 & 1 & Atestação confirmada \cite{anual1932}. \\
1933 & 1 & Atestação confirmada \cite{anual1933}. \\
1934 & 1 & Atestação confirmada \cite{anual1934}. \\
1935 & 1 & Atestação confirmada \cite{ticotico1935}. \\
1936 & 1 & Atestação confirmada \cite{anual1936}. \\
1937 & 1 & Atestação confirmada \cite{anual1937}. \\
1938 & 1 & Atestação confirmada \cite{anual1938}. \\
1939 & 1 & Atestação confirmada \cite{anual1939}. \\
1940 & 1 & Atestação confirmada \cite{anual1940}. \\
1941 & 1 & Atestação confirmada \cite{carioca1941}. \\
1942 & 1 & Atestação confirmada \cite{anual1942}. \\
1943 & 1 & Atestação confirmada \cite{anual1943}. \\
1944 & 1 & Atestação confirmada \cite{anual1944}. \\
1945 & 1 & Atestação confirmada \cite{anual1945}. \\
1946 & 1 & Atestação confirmada \cite{anual1946}. \\
1947 & 1 & Atestação confirmada \cite{jornalnoticias1947}. \\
1948 & 1 & Atestação confirmada \cite{anual1948}. \\
1949 & 1 & Atestação confirmada \cite{anual1949}. \\
1950 & 1 & Atestação confirmada \cite{anual1950}. \\
1951 & 1 & Atestação confirmada \cite{anual1951}. \\
1952 & 1 & Atestação confirmada \cite{anual1952}. \\
1953 & 1 & Atestação confirmada \cite{anual1953}. \\
1954 & 1 & Atestação confirmada \cite{anual1954}. \\
1955 & 1 & Atestação confirmada \cite{camara1955}. \\
1956 & 1 & Atestação confirmada \cite{anual1956}. \\
1957 & 1 & Atestação confirmada \cite{anual1957}. \\
1958 & 1 & Atestação confirmada \cite{anual1958}. \\
1959 & 1 & Atestação confirmada \cite{anual1959}. \\
1960 & 1 & Atestação confirmada \cite{anual1960}. \\
1961 & 1 & Atestação confirmada \cite{anual1961}. \\
1962 & 1 & Atestação confirmada \cite{anual1962}. \\
1963 & 1 & Atestação confirmada \cite{anual1963}. \\
1964 & 1 & Atestação confirmada \cite{anual1964}. \\
1965 & 1 & Atestação confirmada \cite{senado1965}. \\
1966 & 1 & Atestação confirmada \cite{anual1966}. \\
1967 & 1 & Atestação confirmada \cite{anual1967}. \\
1968 & 1 & Atestação confirmada \cite{anual1968}. \\
1969 & 1 & Atestação confirmada \cite{anual1969}. \\
1970 & 1 & Atestação confirmada \cite{anual1970}. \\
1971 & 1 & Atestação confirmada \cite{anual1971}. \\
1972 & 1 & Atestação confirmada \cite{anual1972}. \\
1973 & 1 & Atestação confirmada \cite{anual1973}. \\
1974 & 1 & Atestação confirmada \cite{anual1974}. \\
1975 & 1 & Atestação confirmada \cite{anual1975}. \\
1976 & 1 & Atestação confirmada \cite{anual1976}. \\
1977 & 1 & Atestação confirmada \cite{anual1977}. \\
1978 & 1 & Atestação confirmada \cite{anual1978}. \\
1979 & 1 & Atestação confirmada \cite{correio1979}. \\
1980 & 1 & Atestação confirmada \cite{anual1980}. \\
1981 & 1 & Atestação confirmada \cite{anual1981}. \\
1982 & 1 & Atestação confirmada \cite{anual1982}. \\
1983 & 1 & Atestação confirmada \cite{anual1983}. \\
1984 & 1 & Atestação confirmada \cite{anual1984}. \\
1985 & 1 & Atestação confirmada \cite{camara1985}. \\
1986 & 1 & Atestação confirmada \cite{anual1986}. \\
1987 & 1 & Atestação confirmada \cite{anual1987}. \\
1988 & 1 & Atestação confirmada \cite{anual1988}. \\
1989 & 1 & Atestação confirmada \cite{anual1989}. \\
1990 & 1 & Atestação confirmada \cite{anual1990}. \\
1991 & 1 & Atestação confirmada \cite{anual1991}. \\
1992 & 1 & Atestação confirmada \cite{anual1992}. \\
1993 & 1 & Atestação confirmada \cite{anual1993}. \\
1994 & 1 & Atestação confirmada \cite{folha1994}. \\
1995 & 1 & Atestação confirmada \cite{anual1995}. \\
1996 & 1 & Atestação confirmada \cite{anual1996}. \\
1997 & 1 & Atestação confirmada \cite{anual1997}. \\
1998 & 1 & Atestação confirmada \cite{anual1998}. \\
1999 & 1 & Atestação confirmada \cite{anual1999}. \\
2000 & 1 & Atestação confirmada \cite{anual2000}. \\
2001 & 1 & Atestação confirmada \cite{anual2001}. \\
2002 & 1 & Atestação confirmada \cite{anual2002}. \\
2003 & 1 & Atestação confirmada \cite{anual2003}. \\
2004 & 1 & Atestação confirmada \cite{anual2004}. \\
2005 & 1 & Atestação confirmada \cite{anual2005}. \\
2006 & 1 & Atestação confirmada \cite{anual2006}. \\
2007 & 1 & Atestação confirmada \cite{young2008}. \\
2008 & 1 & Atestação confirmada \cite{anual2008}. \\
2009 & 1 & Atestação confirmada \cite{anual2009}. \\
\bottomrule\end{longtable}\normalsize

\section{Catálogo de pistas fornecidas para ampliação}
Este catálogo preserva as sugestões de jornais, revistas e obras que motivaram a investigação. É uma lista de pistas, não um conjunto de fontes que comprovam as grafias atribuídas. As descrições herdadas são alegações do roteiro; datas de composição, publicação, reedição e disponibilização digital precisam ser distinguidas. A presença de uma obra no catálogo não autentica uma passagem nem valida sua alocação decenal.
\small\begin{longtable}{p{0.09\textwidth}p{0.13\textwidth}p{0.66\textwidth}}
\caption{Pistas herdadas do roteiro, conservadas para auditoria. Descrições não constituem resultados verificados.}\\
\toprule Década alegada & Tipo & Descrição herdada (a verificar) \\\midrule\endfirsthead
\toprule Década alegada & Tipo & Descrição herdada (a verificar) \\\midrule\endhead
1900 & periódico & Gazeta de Notícias (RJ): Nas crônicas satíricas da virada do século, o riso irônico de salão aparecia rigidamente espaçado por pontuações: "Ah! Ah! Ah!". \\
1900 & periódico & O Paiz (RJ): Em artigos de crônica mundana relatando mexericos políticos do início da República. \\
1900 & periódico & Almanak Laemmert (RJ): Nas seções de anedotas e variedades comerciais, registrava o riso clássico pontuado. \\
1900 & obra & Canaã (1902) - Graça Aranha: O clássico do pré-modernismo utiliza a interjeição de forma solene e isolada ao longo dos diálogos. \\
1900 & obra & Recordações do Escrivão Isaías Caminha (1909) - Lima Barreto: Em diálogos que retratam a redação de um jornal carioca da época, os risos são marcados de forma elegante. \\
1900 & obra & A Esfinge (1908) - Coelho Neto: Romance que espelha os salões e teatros cariocas da época da Belle Époque, com diálogos tipograficamente aristocráticos. \\
1910 & periódico & Revista Careta (RJ): Pioneira no uso de piadas curtas ilustradas com a grafia de transição "Hê, hê, hê!". \\
1910 & periódico & Revista O Malho (RJ): O grande semanário satírico ilustrado registrava gargalhadas abertas em charges como "Quá... Quá...". \\
1910 & periódico & A Cigarra (SP): Crônicas paulistas misturando o riso satírico com a modernidade que nascia na capital do café. \\
1910 & obra & Triste Fim de Policarpo Quaresma (1911) - Lima Barreto: O deboche dos burocratas e militares contra o protagonista nacionalista é pontuado com o seco "Ah! ah!". \\
1910 & obra & Urupês (1918) - Monteiro Lobato: O autor introduz o linguajar caipira e regionalista no papel, trazendo a sonoridade crua de risos curtos em suas histórias. \\
1910 & obra & Eu (1912) - Augusto dos Anjos: Embora o livro seja de poesias poéticas e expressionistas, há o registro do riso escarninho da humanidade na literatura da época. \\
1920 & periódico & Correio da Manhã (RJ): Nas crônicas diárias registrando diálogos de salão e fofocas da elite carioca usando "Ah! ah! ah!". \\
1920 & periódico & Diário Nacional (SP): Órgão da imprensa modernista em São Paulo que usava o riso satírico de vanguarda. \\
1920 & periódico & Jornal do Recife (PE): Em seções e crônicas nordestinas. \\
1920 & obra & Macunaíma (1928) - Mário de Andrade: O herói sem nenhum caráter solta risadas folclóricas e descontraídas textualmente ao longo de suas aventuras. \\
1920 & obra & Memórias Sentimentais de João Miramar (1924) - Oswald de Andrade: O uso experimental da linguagem fragmentada reduz o riso a blocos rápidos. \\
1920 & obra & A Menina Morta (Escrito nos anos 20) - Cornélio Pena: Retrata a transição da fala interiorana para o papel com os risos teatrais pontuados. \\
1930 & periódico & Revista Fon-Fon! (RJ): O grande semanário cômico-cultural registrava o deboche de rua com o fonema "Quá! quá! quá!". \\
1930 & periódico & O Radical (RJ): Diário político combativo da época que ironizava adversários políticos fechando notas jornalísticas com escárnios onomatopaicos. \\
1930 & periódico & A Batalha (RJ): Outro veículo forte da imprensa satírica e combativa dos anos 30. \\
1930 & obra & O Quinze (1930) - Rachel de Queiroz: O romance de 30 traz diálogos crus do sertão, transcrevendo risos breves. \\
1930 & obra & São Bernardo (1934) - Graciliano Ramos: O riso cínico ou irônico dos fazendeiros é marcado de forma seca e direta na prosa. \\
1930 & obra & Reforma Agrária (1934) - Marquie: Ensaios e crônicas satíricas da política da década de 1930. \\
1940 & periódico & Diário de Notícias (RJ): Transcrições formais de entrevistas com a rubrica descritiva institucionalizada de (Risos) ou (Gargalhadas). \\
1940 & periódico & A Noite (RJ): Registro das transmissões de rádio ao vivo com rubricas técnicas indicando ao leitor o riso do auditório. \\
1940 & periódico & Folha da Manhã (SP): Registro sóbrio e formal do jornalismo paulistano dos anos 40. \\
1940 & obra & Fogo Morto (1943) - José Lins do Rego: O declínio dos engenhos de açúcar retratado com diálogos fortes e risos de escárnio social. \\
1940 & obra & O Tempo e o Vento: O Continente (1949) - Erico Verissimo: A saga épica gaúcha traz risos de caserna e de taberna descritos literariamente. \\
1940 & obra & Sagarana (1946) - João Guimarães Rosa: O experimentalismo mineiro usa onomatopeias independentes e poéticas para recriar o riso do sertanejo. \\
1950 & periódico & Revista do Pato Donald (1950 - Editora Abril): O berço que massificou o "Quá-quá-quá" com hífens no imaginário jovem brasileiro. \\
1950 & periódico & Tribuna da Imprensa (RJ): Liderado por Carlos Lacerda, o jornal usava o riso sarcástico nas colunas de denúncia política ("He-he-he..."). \\
1950 & periódico & Última Hora (RJ): Jornal moderno e popular de Samuel Wainer, registrando crônicas cotidianas e urbanas. \\
1950 & obra & Grande Sertão: Veredas (1956) - João Guimarães Rosa: Transcreve a fala e o riso dos jagunços através de sonoridades guturais únicas. \\
1950 & obra & Vestido de Noiva (Publicado em livro nos anos 50) - Nelson Rodrigues: O teatro rodriguiano introduz a gargalhada histérica ou maliciosa nas rubricas e falas suburbanas. \\
1950 & obra & Crônicas de Stanislaw Ponte Preta (Iniciadas nos anos 50) - Sérgio Porto: Crônica urbana carioca pré-bossa nova abusando do riso malicioso da noite de Copacabana. \\
1960 & periódico & O Pasquim (1969): O jornal contracultural que revolucionou a linguagem, agrupando letras fluidamente em suas famosas entrevistas descritas ("hahaha" ou "hehehe" soltos). \\
1960 & periódico & Jornal do Brasil (RJ): Com sua famosa reforma gráfica dos anos 60, modernizou a transcrição visual de entrevistas. \\
1960 & periódico & Revista Manchete (RJ): Cobertura da Jovem Guarda e do comportamento jovem usando onomatepeias corridas. \\
1960 & obra & Quarup (1967) - Antônio Callado: O romance da década de 60 traz diálogos descontraídos e informais da juventude e de intelectuais. \\
1960 & obra & A Paixão Segundo G.H. (1964) - Clarice Lispector: O fluxo de consciência foca no sorriso mental ou irônico introspectivo. \\
1960 & obra & Zero (Escrito nos anos 60) - Ignácio de Loyola Brandão: Romance experimental e fragmentado que usa a colagem de textos de rua e onomatopeias pop. \\
1970 & periódico & Gibi da Turma da Mônica (Anos 70 - Editora Abril): Consolidação do "Hi-hi-hi" hifenizado (riso escondido do Cebolinha) e do "Ha-ha-ha". \\
1970 & periódico & Opinião (SP): Jornal semanário de resistência política que usava ironias nas entrelinhas textuais. \\
1970 & periódico & Movimento (SP): Imprensa alternativa dos anos 70 capturando debates com marcações de diálogos. \\
1970 & obra & As Meninas (1973) - Lygia Fagundes Telles: O retrato de três jovens na ditadura usa o riso descontraído de campus universitário em seus diálogos. \\
1970 & obra & Sargento Getúlio (1971) - João Ubaldo Ribeiro: A risada seca e brutal do policial nordestino transcreve a sonoridade sertaneja. \\
1970 & obra & A Morte de D.J. em Paris (1975) - Roberto Drummond: Prosa pop e urbana influenciada diretamente pelas linguagens dos quadrinhos e da publicidade. \\
1980 & periódico & Folha de S.Paulo (SP): Nas colunas sociais cotidianas e bastidores da política da seção Painel, o termo "risos" colapsa e é registrado abreviado como "Rs." ou "Rsrs". \\
1980 & periódico & Jornal do Brasil (RJ): Uso comum do "rs" na coluna de fofocas e cultura informal de Zózimo Barrozo do Amaral. \\
1980 & periódico & O Globo (RJ): Notas rápidas de bastidores da televisão e celebridades adotando a economia mecânica da datilografia. \\
1980 & obra & Viva o Povo Brasileiro (1984) - João Ubaldo Ribeiro: O épico histórico reproduz o riso das variadas classes sociais brasileiras ao longo dos séculos. \\
1980 & obra & O Alquimista (1988) - Paulo Coelho: A literatura de massa dos anos 80 adota uma estrutura limpa e direta de diálogos com interjeições universais de riso ("rires", "sorrisos"). \\
1980 & obra & Feliz Ano Velho (1982) - Marcelo Rubens Paiva: O clássico juvenil dos anos 80 é escrito na exata linguagem coloquial da juventude paulistana da época, com risos corridos e espontâneos de conversa. \\
1990 & periódico & Revista Veja (Anos 90): Nas tradicionais Páginas Amarelas de entrevista, manteve o padrão formal de (Risos) entre parênteses para denotar a fala de políticos e artistas. \\
1990 & periódico & Folhateen (Suplemento jovem da Folha de S.Paulo): O berço impresso que documentava as gírias do mIRC e BBS, registrando graficamente o uso do "rsrsrs" e do "hahaha" que os jovens digitavam nas salas. \\
1990 & periódico & Revista Bizz (Mundo da Música): Jornalismo cultural informal capturando entrevistas com risos transcritos sem rigidez tipográfica. \\
1990 & obra & Chove nos Campos de Cachoeira (1998) - Dalcídio Jurandir: Estilos literários tardios que mantinham viva a tradição descritiva de diálogos no papel. \\
1990 & obra & Budapeste (Escrito/Planejado no fim dos anos 90) - Chico Buarque: Foca na minúcia das palavras e nos dilemas da tradução e grafia de expressões idiomáticas. \\
1990 & obra & Estorvo (1991) - Chico Buarque: Narrativa asfixiante e urbana em que os risos e o humor são tensos, espelhando a transição social dos anos 90. \\
2000 & periódico & Revista Capricho (Anos 2000): A publicação jovem foi uma das primeiras a colocar matérias impressas simulando telas de MSN, imprimindo pela primeira vez termos como "kakakaka" e "kkkkk" em papel. \\
2000 & periódico & Caderno de Informática (O Globo / Folha): Reportagens e colunas da época dedicadas a "traduzir o internetês" dos jovens, publicando tabelas que explicavam o que significavam o kkkk e o vc para os pais assustados. \\
2000 & periódico & Jornal Meia Hora (RJ - Fundado em 2005): Jornalismo popular de massa que quebrou a barreira da sobriedade e passou a adotar manchetes com a linguagem escrachada da internet dos anos 2000. \\
2000 & obra & Os Segredos do Orkut (2005) - Livros de nicho digital: Publicações focadas em desvendar as primeiras grandes redes sociais brasileiras, documentando os logs de depoimentos repletos de "kkkkk". \\
2000 & obra & Manual do MSN Messenger (Meados dos anos 2000): Guias práticos em livrarias ensinando a configurar a ferramenta e a decodificar as abreviações e os novos emoticons rindo. \\
2000 & obra & Literatura de Blog (Coletâneas dos anos 2000): O nascimento de livros impressos baseados em posts de blogs famosos da época (como os textos de Cardosonline ou Melhores do Mundo), que transportaram o "kkkkk" nativo das telas diretamente para o catálogo das editoras brasileiras. \\
\bottomrule\end{longtable}\normalsize

\section{Parecer individual da segunda auditoria}
Os identificadores R01--R66 seguem a mesma ordem do catálogo anterior. A categoria A exige representação do riso autenticada; L identifica referência lexical, temática ou cotejo contextual, sem autenticar a grafia atribuída; B registra identificação bibliográfica; I registra recuperação do exemplar específico sem autenticação da frase; S conserva casos sem ocorrência autenticada ou com identificação insuficiente. Há três A, sete L, dezenove B, quatro I e trinta e três S. Esses estados descrevem o trabalho documental, não a frequência histórica de qualquer grafia. URLs, estado de aquisição e data do parecer constam no CSV público.
\small\begin{longtable}{p{0.07\textwidth}p{0.08\textwidth}p{0.72\textwidth}}
\caption{Parecer individual sobre as 66 indicações. A: representação autenticada; L: referência lexical; B: identificação bibliográfica; I: exemplar específico; S: sem ocorrência autenticada.}\\
\toprule ID & Estado & Resultado e limites \\\midrule\endfirsthead
\toprule ID & Estado & Resultado e limites \\\midrule\endhead
R01 & I & Exemplar de 25/07/1904, n. 207, recuperado; busca no OCR das quatro páginas não autenticou a frase do ovo. OCR degradado impede conclusão negativa definitiva.\\
R02 & S & Busca por título, década e riso não individualizou artigo, edição e página. A descrição genérica não comprova grafia.\\
R03 & S & Sem ano e edição do almanaque ou título da anedota. Busca não individualizou ocorrência; é necessário delimitar volume e seção.\\
R04 & L & Texto integral recuperado: expressão de riso, PDF p. 28. Referência lexical; não autentica o diálogo atribuído ao colono Maciel.\\
R05 & L & Exemplar consultado: segunda edição, 1917, PDF p. 26: Riu-se muito. Não datar a passagem de 1909 sem cotejo da primeira edição.\\
R06 & L & Oito capítulos consultados; capítulo VI contém risos casquinados. É descrição lexical, não uma sequência onomatopaica demonstrada.\\
R07 & S & A busca recuperou Quá em Quase doutor, Careta, 08/05/1915, via Vida urbana. Contexto de recusa, não riso: falso positivo excluído.\\
R08 & I & Edição n. 526 de 12/10/1912 confirmada na capa. Busca nas 58 páginas do OCR não autenticou prefeito/carne barata; desenhos exigem revisão visual integral.\\
R09 & S & Pista alternativa: poema Epístola, A Cigarra, n. 16, 20/01/1915, p. 27, em reprodução externa. Não conferida no exemplar original; não incluída no corpus.\\
R10 & L & Texto de 1915 consultado. Terceira parte, cap. II, contém Você, Quaresma, é um visionario; sem Ah! ah! e sem a continuação sobre plantar. Citação recebida não é transcrição fiel.\\
R11 & A & Quinta edição de 1919: p. impressa 19, PDF p. 24: Quá, quá, quá!, com risos explícitos no contexto. Não retrodatado automaticamente a 1918.\\
R12 & L & Monólogo de uma sombra: Chorando e rindo na ironia infausta. Confirma referência temática, não onomatopeia.\\
R13 & S & Frase de 14/02/1925 não autenticada. Fon-Fon!, maio de 1925, p. 3, noticia retorno do Correio após suspensão; é indício de problema cronológico, não prova autônoma de que não existiu edição em fevereiro.\\
R14 & A & 17/06/1928, n. 290, p. 4: (risos, muito bem), em discurso de Mattos Pimenta. Rubrica descritiva, não grafia de gargalhada.\\
R15 & S & Busca por título, década e risos não individualizou reportagem verificável; faltam dia e página.\\
R16 & A & Edição UFFS de 2019, p. 139: Quá! quá! quá quaquá!... e indicação de gargalhada. Obra original de 1928; a passagem não é a frase Ah! ah! enviada.\\
R17 & L & Exemplar de 1924, p. 114, PDF p. 116: morreu nos lábios de todos o sorriso. Referência lexical; não comprova onomatopeia.\\
R18 & B & BBM-USP registra edição de 1954. A alegação de composição nos anos 1920 não foi demonstrada. Nenhuma passagem de riso autenticada.\\
R19 & I & PDF identificado pelo acervo como n. 23 de 1935 recuperado; capa e OCR consultados. Data de 08/06 e frase sobre rei da Inglaterra não autenticadas; falta conferir expediente legível.\\
R20 & S & Exemplar n. 1175 de 1935 recuperado e OCR pesquisado. Há candidatos lexicais, mas nenhuma grafia atribuída autenticada por imagem.\\
R21 & S & Exemplar n. 31 de 1930 recuperado e OCR pesquisado. Referências lexicais candidatas não autenticam a generalização sobre marcações de riso.\\
R22 & B & Editora Record confirma estreia em 1930; aquisição automática recebeu HTTP 429. Não se confirmou passagem nem grafia.\\
R23 & L & Provas da segunda edição, 1938, recuperadas; PDF p. 63: sorrisos frios. A frase latim de missa não foi localizada no OCR; não retrodatado a 1934.\\
R24 & S & Reforma Agrária (1934), de Marquie: buscas por título, ano e variante do nome não identificaram registro bibliográfico suficiente. Não tratar como obra inexistente.\\
R25 & I & Edição de 18/03/1944, n. 6564, confirmada no cabeçalho. Busca do OCR não autenticou ministro/boatos de racionamento; texto de baixa qualidade.\\
R26 & S & Busca por título, ano e rubrica não individualizou transmissão ou reportagem com página. Não demonstrada a descrição recebida.\\
R27 & S & Busca produziu falso positivo: Jornal de Notícias, de 02/03/1947, e não Folha da Manhã. Documento incluído separadamente; não substitui esta referência.\\
R28 & B & Bibliografia da Academia Brasileira de Letras confirma Fogo morto, 1943. Texto e passagem sobre riso não conferidos.\\
R29 & B & Editora identifica O Continente na trilogia publicada entre 1949 e 1962. Metadado não autentica diálogo de caserna ou taberna.\\
R30 & B & Museu Casa Guimarães Rosa registra exemplar de Sagarana de 1946. Frase He! he! e mundo dá muitas voltas não autenticada.\\
R31 & S & Revista de 1950 e alegado Almanaque de 1953 não são a mesma referência. Faltam número, história e página para verificar moedas e Quá-quá-quá.\\
R32 & S & Busca de frase e jornal não localizou plano é honesto. Ano de 1955 sem dia, coluna e página não individualiza o documento.\\
R33 & S & Busca por jornal, década e risos não individualizou crônica. É necessária referência de edição.\\
R34 & B & Página editorial de Grande sertão: veredas consultada; primeira edição em 1956. Nenhuma passagem específica de riso foi autenticada.\\
R35 & B & Registro institucional sobre Vestido de noiva identifica escrita em 1943; estreia e publicação são eventos distintos. A edição em livro dos anos 1950 não foi individualizada.\\
R36 & S & Crônicas de Stanislaw Ponte Preta não designa título, edição ou crônica específica. Identidade do cronista não confirma data e grafia de passagem.\\
R37 & S & Coleção identificada no dossiê da Biblioteca Nacional. URL tentada para n. 12 retornou 404; diálogo atribuído a Jaguar permanece não autenticado.\\
R38 & S & A reforma gráfica do jornal não demonstra uso de hahaha. Busca não localizou entrevista e página específicas na década de 1960.\\
R39 & S & Sem edição e entrevistado identificados. Busca por título, década e rubrica não autenticou onomatopeia corrida.\\
R40 & B & Página editorial de Quarup localizada; download automático HTTP 429. Frase sobre jornais não autenticada e sem página.\\
R41 & B & Registro museológico localizado para A paixão segundo G.H.; download retorna página sem conteúdo textual suficiente. Não confirma sorriso mental ou grafia.\\
R42 & B & ABL distingue edição italiana de 1974 e portuguesa de 1975. Alegada composição nos anos 1960 não autoriza contar publicação naquela década.\\
R43 & S & Menção genérica a gibis dos anos 1970. Busca não individualizou número e história com as formas atribuídas.\\
R44 & S & Busca por Opinião, década e risos não individualizou passagem verificável. A caracterização ideológica não prova grafia.\\
R45 & S & Busca por Movimento, década e risos não individualizou debate com localizador. Pendente identificação de edição.\\
R46 & B & Página editorial de As meninas consultada. Frase sobre aposta e Ha-ha-ha! não autenticada; sem acesso integral conferido.\\
R47 & B & Página editorial localizada na busca; aquisição automática falhou com 404. Data de obra não é prova de uma risada específica.\\
R48 & B & Coletânea de 1975 identificada também em pesquisa institucional UFRN; catálogo WorldCat bloqueou aquisição com 403. Nenhuma passagem de riso confirmada.\\
R49 & S & Busca da data 14/11/1986 e frase de congelamento não autenticou Rs. na seção Painel. Rubrica risos de outra publicação de 1994 não substitui esta alegação.\\
R50 & S & Busca por Zózimo, década e rs não individualizou nota. Não demonstrado uso comum de rs no jornal nos anos 1980.\\
R51 & S & Busca por jornal, década e rsrs não individualizou nota. Alegada economia da datilografia permanece explicação sem evidência específica.\\
R52 & B & Página editorial de Viva o povo brasileiro consultada. Descrição literária genérica não comprova forma particular de riso.\\
R53 & B & Página editorial de O Alquimista consultada. Sorriso e rir são itens lexicais; não devem ser classificados automaticamente como interjeições.\\
R54 & B & Página editorial de Feliz ano velho consultada. Frase burradas do colégio e sequência Ha ha ha não autenticadas.\\
R55 & S & Busca de Veja, 07/05/1993 e frase não fugirei da raia não individualizou entrevista. Data e entrevistado continuam sem autenticação.\\
R56 & S & Busca de Folhateen, junho de 1997, mIRC e grafias não identificou reportagem com dia e página. Não confirmada a frase recebida.\\
R57 & S & Busca de Bizz, década e hahaha não identificou entrevista específica. Sem edição, autoria e página.\\
R58 & B & Fonte institucional UEPA registra primeira publicação em 1941. Ano de 1998 pode designar edição; não é data de nascimento da obra nem da passagem.\\
R59 & B & Página editorial informa lançamento em 10/09/2003. Planejamento nos anos 1990 não demonstrado e não equivale a publicação.\\
R60 & B & Página editorial informa lançamento em 02/08/1991. Nenhuma passagem sobre riso foi conferida no texto integral.\\
R61 & S & Busca de Capricho, setembro de 2007 e frase quase quebrou a cara não individualizou reportagem. Não se sustenta afirmação de primeira impressão de kkkk.\\
R62 & S & Caderno de Informática de dois jornais é referência ambígua. Busca por internetês e kkkk não individualizou tabela de abreviações com data e página.\\
R63 & S & Busca por Meia Hora, 2005 e kkkk não identificou manchete. Fundação do veículo não demonstra adoção de grafia específica.\\
R64 & S & Título, autor e editora do alegado livro de 2005 não identificados. Resultado homônimo é postagem de blog de 2008, que não autentica o livro.\\
R65 & S & Manual do MSN Messenger sem autor, ISBN ou editora. Busca não individualizou obra e edição; não autentica convenções descritas.\\
R66 & S & Literatura de Blog é categoria, não referência bibliográfica individual. Cardosonline e Melhores do Mundo não podem ser convertidos em livro sem título, autor e edição.\\
\bottomrule\end{longtable}\normalsize

\section{Correções conceituais e decisões editoriais}
\begin{enumerate}[leftmargin=*]
\item Mudança histórica e motivação icônica são dimensões distintas, que podem coexistir.
\item Continuidade lexical não exige afixação produtiva ou radical superficialmente preservado.
\item Uma cadeia histórica pode ser reconstruída sem documentação de todas as gerações. A ausência de elos observados reduz a força da demonstração, mas não prova independência.
\item Não há prova de que \forma{quiá} desapareceu da circulação social ao ser impresso. A extensão de uso requer corpus.
\item \forma{Quiá} em \textit{Til} não é ocorrência literal de \forma{kkkk}, nem demonstra a primeira representação de riso da língua.
\item \forma{Você} para \forma{cê} envolve aférese. A cronologia que situa \forma{você} apenas no século XX contradiz a atestação de 1638 discutida por Monte.
\item \forma{Cá} é nome de letra ou sílaba, não um fonema único; /b/ é bilabial; a gargalhada não contém necessariamente /k/.
\item \forma{Vilão} não se explica como simples mudança de \forma{-us} em \forma{-ão}. Mudança semântica não determina sozinha classe gramatical.
\item Famílias linguísticas diferentes não excluem empréstimo; \forma{chikutaku} não é prova de isolamento.
\item Datas de plataformas, reformas ortográficas e teclados são contexto potencial, não provas de origem de uma sequência.
\item Preferências geracionais não são universais; o corpus não estima usos brasileiros da geração Alpha.
\item Uma referência de divulgação citada num artigo acadêmico não se torna novo testemunho primário.
\item A quota de três jornais e três livros é editorial. A justificativa estatística depende do estimando e da coleta.
\item Datas precisas e aspas exigem localizador e conferência. Ilustrações não podem figurar como transcrições.
\end{enumerate}

\section{Dicionário de dados e reprodução}
\texttt{attestations.csv} reúne testemunhos; \texttt{claims\_audit.csv}, os 22 alvos; \texttt{annual\_coverage.csv}, a cobertura anual; \texttt{candidate\_sources.csv}, as pistas; \texttt{downloads.json}, as aquisições; \texttt{detection\_sensitivity.csv}, os cenários matemáticos; \texttt{summary.json}, as contagens.

\texttt{build\_data.py} gera tabelas de registros codificados e parâmetros declarados. A aquisição é separada da análise para que falhas de rede não sejam interpretadas como resultados linguísticos. Regex localizam candidatos e não substituem anotação. A segunda rodada acrescenta os arquivos \texttt{references\_audit.csv}, \texttt{claims\_verification.csv} e \texttt{extra\_attestations.json}, além dos registros de buscas. O repositório documenta comandos e versões.

Não foram produzidos logs sintéticos para os anos 1990, participantes fictícios, contagens populacionais ou intervalos de confiança para a amostra dirigida. Os cenários de planejamento são identificados como hipotéticos. A distinção entre realizado e proposto integra a metodologia.
